% Created 2026-06-16 ti 10:43
% Intended LaTeX compiler: pdflatex
\documentclass[11pt]{article}

\usepackage[utf8]{inputenc}
\usepackage[T1]{fontenc}
\usepackage{graphicx}
\usepackage{longtable}
\usepackage{wrapfig}
\usepackage{rotating}
\usepackage[normalem]{ulem}
\usepackage{amsmath}
\usepackage{amssymb}
\usepackage{capt-of}
\usepackage{hyperref}
\usepackage{listings}
\usepackage{color}
\usepackage{amsmath}
\usepackage{array}
\usepackage[T1]{fontenc}
\usepackage{natbib}
\date{\today}
\title{}
\begin{document}

\tableofcontents

\section{Notes}
\label{sec:orgd88496e}

\subsection{Input for simulate cohort}
\label{sec:orgbf1423a}
p contains,
\begin{itemize}
\item Maximum number of events for the sim.
\item How baseline parameter values are distributed
\item How events/visits are distributed.
\item parameter values:
(intercept, relating to expected values of parameters and the
timeline of events)
(scale, relating to likelihood of events)
(effect, effects of different parameters on eachother,
effects of different parameters/events on likelihood of events,
as well as treatments on the parameters/events)
\end{itemize}
Effects.
Variables --> Variables
Variables --> Events (intermediate/absorbing)
Intermediate --> Events
No
Events --> Variables? Why should measurements in general not change
as a result of an events?
\subsection{What happens to simulated data for different values of scale\textsubscript{dropout}}
\label{sec:org4c0d2a0}
scale\textsubscript{dropout} < 0 (No other events past baseline)
scale\textsubscript{dropout} = 0 (In general goes to the maximum amount of events,
max\textsubscript{follow} = 6, but dropouts still occur for most)
scale\textsubscript{dropout} > 0 (Increasing rate of dropout events)
\section{{\bfseries\sffamily DONE} Randomized treatment Diabetes data generation (similar time-series to Kathrine data)}
\label{sec:org8c759f1}
\begin{itemize}
\item Create setting p (distributions/schedule/parameters/effects) that
will mimick time-series of LEADER study data, through simulate
cohort. (EXAMPLE SETTING DONE). First try shared with Thomas.
\end{itemize}
\section{{\bfseries\sffamily DONE} Need to modify to handle variance different from one apparently.}
\label{sec:org5b6b2e9}
\section{{\bfseries\sffamily TODO} Visits do not seem to follow the schedule. Check it out.}
\label{sec:orgf234d6a}
Okay seems to me that the scheduling in some sense "fails" - at least
if the intention is to have visits every 6 months. However they are
instead spaced such that the visit happens exactly 6 months after the
previous event. Is this the intention? ASK thomas.

Also weird that visits "compete" with events\ldots{}
Try fix? Remember to mark changes in functionality.

Should the biomarkers be updated at every type of event?

Have made baseline treatment (either lira or placebo) randomized.
\subsection{Talk with Kathrine at some point}
\label{sec:org688313b}
\begin{itemize}
\item Have rough outline from looking at the leader targets pipeline.
Have made assumptions about distributions of differnet variables.

\item I think the only measurement variables listed in the pipeline
is a1c level/change, which i think works nicely in the simulation.
Should the rest of the baseline biomarkers be updated as well?

\item Then generally what effects between covariates/events are relevant.

\item Is there a large difference in the likelihood between events
\end{itemize}

What do i need specifically?
\begin{itemize}
\item Mean/variance values that make sense.
\item Scales of event probabilities (Maybe relative is enough)
\item Changes in distributions?
\item Update biomarkers?
\end{itemize}
\section{{\bfseries\sffamily TODO} Simulate Diabetes register type data}
\label{sec:orgaa5c950}
\begin{itemize}
\item create setting with appropriate variables/distributions/events etc
making timeseries that mimicks register data.
\end{itemize}
\subsection{Overview}
\label{sec:org6dec8d9}
\begin{itemize}
\item First step is to recreate simulated timeseries of correct
names/variables. Trying one register at a time (a lot of events
may be NULL).
Ok so thing is, if they are independently simulated how does
one match ids from different registers. Does not really make
sense to do it like that i suppose. Make together (entire
history of person is followed) such that the progression of
the registers is alligned - Then split registers (if made for
test of get\textsubscript{DMreg}()).
\item What makes sense to make here. Should it be some population
with some percentage chance of getting a diabetes entry.
\end{itemize}
\subsection{How does the typical history look?}
\label{sec:orgcfaa026}
Most likely they get diagnosis/lab first then prescription. No foot first
right? What about the others. Maybe look at 3740 for shape and overview
of the data.

Need data like: MFR, PDB, LPR, LMDB, LABKA, (SSR, DVDD)
\section{Other thing causal changes in treatment}
\label{sec:org858473b}
\end{document}
